\section{Typical: the decision pass}
\label{sec:method}

Typical turns a pretrained causal language model into a decision function whose output is read from hidden
states in one pass over a short suffix.

\subsection{Typed decisions}
\label{sec:method:task}

\begin{definition}[Typed decision]
\label{def:task}
A decision is a tuple $\decision = (\state, \query, \cands, \dtype)$: a state $\state$ (the document,
transcript or trace the decision is about), a question $\query$ (the instruction or rubric), a candidate
set $\cands = \{\cand{1},\dots,\cand{\numcands}\}$ of label strings supplied at request time, and a type
$\dtype \in \{\textsc{Choice}, \textsc{Score}, \textsc{Noul}\}$. The model returns
\[
  P(\cand{j} \mid \state, \query, \cands),\ j = 1,\dots,\numcands,
  \qquad
  \pnull = P(\nullsym \mid \state, \query, \cands),
  \qquad
  \pnull + \textstyle\sum_{j} P(\cand{j} \mid \state, \query, \cands) = 1,
\]
and emits no token.
\end{definition}

The type is a contract on the returned distribution, enforced by the head and the training target.
\emph{Choice} is categorical over $\cands$. \emph{Score} is categorical over ordered levels, where near
misses cost less than distant ones. \emph{Noul} is a yes/no decision: a Bernoulli probability for one proposition, with
$\cands = \{\text{yes}, \text{no}\}$, invariant to the order in which the caller lists the two labels. The
abstention event $\nullsym$ is available under every type and is never one of the candidates. Because
$\cands$ is text supplied per call, the model has no output index for any label and must read the
candidates to score them.

\subsection{Cached state and causal suffix}
\label{sec:method:pass}

Typical encodes the state once as an ordinary causal prefix and retains its key-value cache. Each question
against it is a short causal suffix: the question text, the candidates rendered as lettered lines, an
answer cue, and a terminal decision token. The suffix attends to the cached state, and a batch of
suffixes against one state runs in chunks without attending to one another. A decision is therefore a
forward pass on $\mathrm{concat}(\state, \query, \cands)$ with the state half read from cache, which an
equivalence test on the real model confirms (\Cref{app:cache-equivalence}).

\subsection{The contextual-candidate readout}
\label{sec:method:readout}

At the tap layer the suffix pass yields the terminal decision state $\hD$, the hidden state at the
decision token, and for each candidate a contextual state $\hctx{j}$, the mean hidden state over the
tokens of candidate $j$'s rendered span. Candidate $j$ receives the score
\begin{equation}
  u = W_h \hD, \qquad v_j = W_c\, \hctx{j}, \qquad
  s_j = \frac{u^{\top} v_j}{\sqrt{d_v}} + w^{\top}(u \odot v_j) + g(u \odot v_j),
  \label{eq:readout}
\end{equation}
where $W_h$ and $W_c$ are layer-normalized linear projections to width $d_v$, $w$ is a vector and $g$ a
two-layer MLP: a bilinear form with a small learned correction on the elementwise product. The option
letters stay in the rendered text and are never read out. Because the candidates are tokens in the same
causal computation as the state and the question, $\hctx{j}$ carries the candidate's meaning together
with its slot and its set, and $\hD$ has attended to all of them.

\Cref{sec:analysis:readout} compares three readouts. \emph{N1}, the letter-logit readout, scores $\hD$
against the output embeddings of the option letters (their next-token logits, with a softmax null),
the standard multiple-choice interface read without generating. \emph{N2}, the candidate-blind semantic readout, applies \Cref{eq:readout} to an embedding of
the candidate text from a frozen text-embedding model, which carries meaning and no slot identity.
\emph{N3}, the contextual-candidate readout Typical deploys, applies \Cref{eq:readout} to $\hctx{j}$.
N2 and N3 share the scorer, the factored null and the feature normalization and differ only in the
candidate vector. N1 and N3 differ in the readout and in the null form native to each; option shuffling
during training is shared by all three.

\subsection{Mid-depth tap}
\label{sec:method:depth}

We truncate the backbone at about 71\% of its depth: the readout taps layer 20 of 28 at 1.7B, 26 of 36 at
4B and 8B, and 28 of 40 at 14B, and the layers above the tap are never computed. LoRA \citep{hu2022lora}
at rank 16 adapts the top eight kept layers, and everything below is frozen. We chose the tap for what
the representation holds: our first full-scale model read the final layer and scored 58.6 on SNLI, the
hypothesis-only baseline, because it ignored the premise. The mid-depth tap carries evidence and
question-conditioned knowledge in one model (\Cref{sec:analysis:depth}), and every released and ladder
model uses it.

\subsection{Typed heads}
\label{sec:method:typed}

The three types share the backbone and the tap and differ in the training target and, for Noul, the
terminal head. \emph{Choice} is the $\numcands$-way readout of \Cref{eq:readout}, trained with
cross-entropy against the gold distribution over $\cands \cup \{\nullsym\}$ through the gate of
\Cref{sec:method:null}. \emph{Score} uses the Choice readout over the rendered levels with an
ordinal-smoothed target $\tilde{t}_j \propto \exp(-|j - y|/\tauord)$ for gold level $y$, normalized over
the valid levels, with $\tauord = 0.7$; $\tauord \to 0$ recovers the one-hot Choice target, and the change
adds no parameters. \emph{Noul} uses a dedicated Bernoulli head, $P(\text{yes}) = \sigma(w_n^{\top}
\hD)$, on a query-only suffix that renders no candidate text, with the complement assigned to ``no''.
Abstention on this path uses a gate of the form of \Cref{eq:factored-null} with its own parameters, whose
score statistics come from $\log P(\text{yes})$ and $\log P(\text{no})$ and whose state input is $\hD$,
since no candidate span is rendered. Every gate input is symmetric in the two labels, so for rows routed
to the Bernoulli head the returned yes, no and $\nullsym$ probabilities are exactly invariant to label
order.

Routing is per row, because a batch mixes types. A row takes the Bernoulli path if and only if its
candidate set is exactly $\{\text{yes}, \text{no}\}$, Score rows take the ordinal target, and every other
row is scored by the Choice path, bit-identical to a model with the typed heads disabled
(\Cref{app:typed-primitives}).

\subsection{Factored abstention gate}
\label{sec:method:null}

Abstention is a separate gate outside the candidate softmax. Given the scores $s_j$ of
\Cref{eq:readout}, the gate reads a vector $z$ of permutation-invariant statistics over the valid
candidates (maximum score, top-2 margin, mean score, log-mean-exp of the scores), the projected decision
state, and the mean and variance of the projected candidate states $v_j$, and computes
$r = \sigma(g_\nullsym(z))$. The returned distribution is
\begin{equation}
  \label{eq:factored-null}
  P(\cand{j} \mid \state, \query, \cands) = (1-r)\,\frac{\exp(s_j/T)}{\sum_{k=1}^{\numcands}\exp(s_k/T)},
  \qquad
  \pnull = r,
\end{equation}
implemented as logits $\ell_j = \log \mathrm{softmax}(s/T)_j + \log(1-r)$ and $\ell_\nullsym = \log r$,
so one cross-entropy over $\numcands + 1$ outcomes trains both parts. Two properties follow exactly.
(i)~\emph{Pairwise non-null odds are invariant to abstention}: for fixed scores
$P(\cand{j})/P(\cand{k}) = \exp((s_j - s_k)/T)$, independent of $r$, because $(1-r)$ multiplies every
non-null candidate; the contextual scores themselves remain set-dependent. (ii)~\emph{The gate is
temperature-invariant}: $T$ enters only the candidate softmax and $r$ reads untempered statistics, so
recalibrating $T$ never moves $\pnull$. The gate is symmetric in the pairs $(s_j, v_j)$, so it adds
no order dependence, while $r$ depends on what was offered by design. We never render the null as a candidate line, because a rendered line competes
for softmax mass that dilutes as $\numcands$ grows; \Cref{sec:analysis:null} removes that line from a model
that already has the gate and measures the difference.

\subsection{Cost model}
\label{sec:method:cost}

$M$ decisions against one state cost $O(|\state|) + M \cdot O(|\query| + |\cands|)$ token positions, where
$|\cands|$ counts the rendered candidate tokens, and every position runs through only the kept 71\% of the
layers. The state term is paid once and amortized over every question asked of it; the per-decision term
grows with the question and the candidate text and is independent of the other questions. In our research harness on one H100
at $\numcands = 2$, a decision takes 45, 56 and 60\,ms at 1.7B, 4B and 14B, so an $8\times$ larger backbone
is $1.3\times$ slower. At large $\numcands$ the rendered candidates dominate the per-decision term, and in
throughput mode a candidate-blind scorer, whose cost does not grow with the rendered set, is cheaper per
query by a margin that grows with $\numcands$ (\Cref{sec:results:cost}).
